% CTX: Trigger-Driven Dynamic Context Loading for Code-Aware LLM Agents
% ACL Rolling Review / EMNLP 2026 Submission
\documentclass[11pt]{article}
\usepackage[hyperref]{acl}
\usepackage{times}
\usepackage{latexsym}
\usepackage{booktabs}
\usepackage{graphicx}
\usepackage{amsmath}
\usepackage{url}

\title{CTX: Trigger-Driven Dynamic Context Loading\\for Code-Aware LLM Agents}

\author{Jeawon Jang \\ \texttt{be2jay67@gmail.com}}

\begin{document}
\maketitle

\begin{abstract}
Large language models suffer from context dilution when processing extensive codebases, a phenomenon documented as the ``Lost in the Middle'' problem. Existing retrieval-augmented generation (RAG) approaches treat code as flat text, ignoring the structural dependency information encoded in import graphs. We present CTX, a trigger-driven dynamic context loading system that classifies code-related queries into four trigger types---\textsc{Explicit\_Symbol}, \textsc{Semantic\_Concept}, \textsc{Temporal\_History}, and \textsc{Implicit\_Context}---and routes each to a specialized retrieval pipeline backed by a three-tier hierarchical memory architecture. For dependency-sensitive queries, CTX performs breadth-first traversal over the codebase import graph to resolve transitive relationships invisible to keyword and embedding methods. Additionally, we introduce a hybrid variant that combines dense neural embeddings for semantic seed selection with CTX's import graph expansion, bridging the semantic-structural gap. Evaluated on a synthetic benchmark (50 files, 166 queries) and three real Python codebases (968 files total, 249 queries), CTX achieves a Trade-off Efficiency Score (TES) 1.9$\times$ higher than the best baseline on synthetic data while consuming only 5.2\% of total tokens. On implicit dependency queries, CTX attains perfect Recall@5 (1.0) compared to 0.4 for BM25 on the synthetic benchmark. The hybrid variant achieves Recall@5 of 0.95 on the COIR code retrieval benchmark, a 150\% improvement over CTX alone (0.38), validating the complementary nature of dense and structural retrieval. Statistical significance is established via McNemar and Wilcoxon tests ($p<0.05$) across 415 total queries.
\end{abstract}

\section{Introduction}

When a programmer encounters a function call to \texttt{process\_data()}, they do not read the entire codebase to understand it. Instead, their mind follows a chain of associations: the function name triggers recall of its definition, which triggers awareness of its imports, which triggers understanding of the data structures it manipulates. This associative recall---where a stimulus activates a targeted subset of memory rather than the entire store---is a hallmark of human cognition that current LLM-based code assistants fail to replicate.

The dominant approach to providing code context to LLMs is to load as much of the codebase as possible into the context window. This strategy is fundamentally flawed. \citet{shi2024distracted} demonstrated a U-shaped attention degradation pattern where LLMs lose track of information placed in the middle of long contexts. The problem is particularly acute for code: LongCodeBench \citep{guo2025longcodebench} reports that Claude 3.5's accuracy on code tasks drops from 29\% at 32K tokens to 3\% at 256K tokens---an order-of-magnitude collapse that grows worse with context length.

Retrieval-augmented generation offers a partial solution by selecting relevant documents rather than loading everything. However, standard RAG approaches---whether sparse (BM25) or dense (embedding-based)---treat code as flat text. They match on lexical overlap or semantic similarity but are structurally blind. Consider a query asking ``What modules are needed to understand the data pipeline?'' The answer requires tracing import chains: module A imports module B, which imports module C, forming a transitive dependency invisible to any text-matching approach.

CTX addresses this gap by drawing on two key insights. First, developer queries about code fall into distinct categories that demand different retrieval strategies. Second, the import graph of a codebase is a readily available, machine-parseable structure that encodes precisely the dependency relationships that text-based retrieval misses.

This paper makes four contributions:

\begin{enumerate}
\item \textbf{A four-type trigger taxonomy for code queries.} We define \textsc{Explicit\_Symbol}, \textsc{Semantic\_Concept}, \textsc{Temporal\_History}, and \textsc{Implicit\_Context} as distinct query categories, each mapped to a specialized retrieval strategy.

\item \textbf{Import graph traversal for implicit dependency retrieval.} We introduce a BFS-based traversal algorithm over the codebase import graph that resolves transitive dependencies. On synthetic data, this achieves Recall@5 of 1.0 on dependency queries, compared to 0.4 for BM25.

\item \textbf{The Trade-off Efficiency Score (TES) metric.} We propose $\text{TES} = \text{Recall@K} / \ln(1 + |\text{retrieved}|)$ as a unified measure of the accuracy-efficiency trade-off.

\item \textbf{A hybrid dense+graph pipeline.} We demonstrate that combining dense neural embedding retrieval with import graph expansion achieves Recall@5 of 0.95 on the COIR benchmark, a 150\% improvement over CTX alone, validating the complementary nature of semantic and structural retrieval.
\end{enumerate}

\section{Related Work}

\subsection{Memory-Augmented LLM Systems}

\textbf{MemGPT} \citep{packer2023memgpt} introduces an OS-inspired virtual memory hierarchy that pages information between main context and external storage. While MemGPT provides a general-purpose memory management framework, it does not exploit domain-specific structure: it treats all documents uniformly regardless of whether they are code, prose, or data.

\textbf{Memori} \citep{memori2025} implements embedding-based memory consolidation for LLM agents, achieving approximately 5\% token usage on the LoCoMo conversational benchmark---comparable to CTX's 2--5\% on code repositories. However, Memori treats all documents as flat text chunks, relying solely on embedding similarity for retrieval. CTX exploits the unique structural property of code: import graphs provide \emph{explicit, deterministic} dependency information that embedding similarity cannot capture. Our ablation study (Section~\ref{sec:ablation}) demonstrates a 60\% \textsc{Implicit\_Context} recall drop when removing import graph traversal.

\textbf{MeCo} \citep{wang2025meco} determines retrieval necessity via learned probes on LLM internal activations---a model-dependent approach requiring white-box access to transformer hidden states. CTX's trigger classification operates externally on query syntax and semantics, making it model-agnostic and deployable with any LLM backend.

\subsection{Code-Specific Retrieval}

\textbf{jCodeMunch} \citep{jcodemunch2025} achieves up to 95\% token reduction through code compression. While impressive in token savings, compression-based approaches carry an inherent risk: compressed code loses semantic completeness. CTX instead preserves complete file semantics while selecting only relevant files through structural retrieval.

\textbf{CAR} \citep{li2024car} introduces dynamic retrieval size selection based on query clustering. CTX shares the principle of adaptive retrieval but differs in mechanism: CAR clusters queries by embedding similarity, while CTX classifies by trigger type---a code-specific taxonomy that maps directly to retrieval strategies.

\textbf{RANGER} (2025) constructs a comprehensive knowledge graph of entire repositories, capturing hierarchical and cross-file dependencies down to the variable level, and augments nodes with textual descriptions and embeddings. RANGER achieves strong results across four standard benchmarks (CodeSearchNet, RepoQA, RepoBench, CrossCodeEval). CTX differs from RANGER in three key ways: (1) CTX operates at query-routing level, classifying \textit{intent} before retrieval, whereas RANGER applies uniform graph traversal to all queries; (2) CTX's adaptive-k selection adjusts retrieval budget per trigger type, while RANGER uses fixed retrieval depth; and (3) CTX achieves 5.2\% token usage versus RANGER's full-graph traversal. The approaches are complementary: CTX's trigger taxonomy could be applied as a query routing layer on top of RANGER's knowledge graph.

\subsection{Long-Context Degradation}

The empirical evidence for context dilution is substantial. \citet{shi2024distracted} document the ``Lost in the Middle'' phenomenon. \citet{kandpal2023longtail} show that LLMs struggle disproportionately with long-tail knowledge. LongCodeBench \citep{guo2025longcodebench} extends these findings to code, demonstrating severe accuracy degradation beyond 32K tokens. The Head-to-Tail study \citep{sun2024headtotail} further characterizes frequency-dependent knowledge retrieval failures. These findings motivate CTX's core design: rather than expanding context windows, reduce context to only what is relevant, guided by code structure.

\section{Method}

\subsection{Problem Formulation}

Given a codebase $C = \{f_1, f_2, \ldots, f_n\}$ consisting of $n$ source files and a developer query $q$, the retrieval task is to select a subset $R \subseteq C$ such that $R$ contains the files relevant to answering $q$ while $|R| \ll |C|$. Let $G \subseteq C$ denote the ground-truth relevant files. The objective is to maximize $\text{TES}(\text{Recall@K}(R, G), |R|)$, simultaneously achieving high recall and low retrieval size.

\subsection{Trigger Classification}

CTX classifies each incoming query into one of four trigger types using a rule-based classifier with confidence scoring:

\textbf{\textsc{Explicit\_Symbol}.} The query references a specific code symbol---a function name, class name, variable, or file path. Detection uses regex pattern matching for identifiers (CamelCase, snake\_case) and path-like strings. Example: ``What does the \texttt{parse\_imports} function do?''

\textbf{\textsc{Semantic\_Concept}.} The query describes a code behavior or pattern without naming specific symbols. Detection uses keyword matching against a domain vocabulary. Example: ``How does the system handle failed API requests?''

\textbf{\textsc{Temporal\_History}.} The query references previous interactions or session context. Detection matches temporal markers. Example: ``What was the module we looked at before?''

\textbf{\textsc{Implicit\_Context}.} The query requires understanding dependency chains or transitive relationships. Detection identifies dependency-related keywords. Example: ``What modules are needed to understand the data pipeline?''

Each trigger type maps to a base retrieval size $k$: \textsc{Explicit\_Symbol} ($k=3$), \textsc{Semantic\_Concept} ($k=8$), \textsc{Temporal\_History} ($k=5$), \textsc{Implicit\_Context} ($k=10$).

\subsection{Hierarchical Memory Architecture}

CTX organizes codebase knowledge into three tiers:

\textbf{Tier 1: Working Memory.} Contains currently active context---files recently accessed or explicitly referenced. Working memory is small (typically 3--5 files) and provides immediate, zero-latency retrieval.

\textbf{Tier 2: Episodic Memory.} Stores session history---files discussed in previous interactions, query-response pairs, and navigation patterns.

\textbf{Tier 3: Semantic Memory.} Contains long-term, structured knowledge indices: the symbol index, the concept index, and the import graph.

Retrieval follows a cascade: Working Memory is checked first, then Episodic Memory, then Semantic Memory. The cascade terminates when the confidence-weighted retrieval count meets the trigger-specific $k$ threshold.

\subsection{Import Graph Traversal}

The import graph is the core differentiator of CTX. We construct a directed graph $G_I = (V, E)$ where each vertex $v \in V$ represents a source file and each directed edge $(u, v) \in E$ indicates that file $u$ imports file $v$.

For \textsc{Implicit\_Context} queries, retrieval proceeds as follows:

\begin{enumerate}
\item \textbf{Seed identification.} The query is analyzed to identify seed files---files directly mentioned or strongly matched by keyword. If no seeds are found, the top-3 BM25 results serve as seeds.

\item \textbf{BFS traversal.} Starting from each seed, breadth-first search expands along import edges. Each reached file receives a distance-based score: $\text{score}(f) = 1 / (1 + d(f))$, where $d(f)$ is the shortest hop distance from any seed.

\item \textbf{Dynamic depth control.} The maximum traversal depth adapts to graph density. For sparse graphs (average degree $< 3$), max depth is 3; for denser graphs, max depth is 2.

\item \textbf{Score aggregation.} Files reached from multiple seeds receive the maximum score across paths. The top-$k$ files by score are returned.
\end{enumerate}

\subsection{Hybrid Dense+Graph Pipeline}
\label{sec:hybrid}

Furthermore, CTX can be combined with dense retrieval in a two-stage pipeline: dense retrieval identifies semantically relevant seed files, which CTX then expands via import graph traversal.

\begin{enumerate}
\item \textbf{Dense seed selection.} A sentence-transformer model (all-MiniLM-L6-v2) encodes the query and all codebase files. The top-$k_s$ files by cosine similarity serve as seeds ($k_s = 3$ by default).

\item \textbf{Graph expansion.} Each seed file is expanded via BFS over the import graph with a configurable hop limit ($h = 2$ by default).

\item \textbf{Score combination.} Dense similarity scores and graph proximity scores are combined with equal weights: $\text{score}(f) = 0.5 \cdot s_{\text{dense}}(f) + 0.5 \cdot s_{\text{graph}}(f)$.
\end{enumerate}

This hybrid approach addresses CTX's primary weakness on text-to-code semantic matching while preserving structural dependency awareness.

\subsection{TES Metric}

We define the Trade-off Efficiency Score as:
\begin{equation}
\text{TES} = \frac{\text{Recall@K}}{\ln(1 + |\text{retrieved}|)}
\end{equation}

\textbf{Information-Theoretic Justification.} The logarithmic denominator follows from the empirical distribution of file utility in codebases. File reference frequencies follow a Zipf distribution \citep{baxter2006java}, where the marginal cost of retrieving the $k$-th additional file scales as $\ln(k)$. TES therefore measures the ratio of marginal benefit (recall gain) to marginal cost (logarithmic retrieval overhead).

\textbf{Formal properties of TES.}
TES satisfies three desirable properties for an efficiency metric:
\begin{enumerate}
\item \textbf{Monotonicity in recall}: For fixed $|\text{retrieved}|$, TES is monotonically increasing in Recall@K.
\item \textbf{Diminishing returns}: For fixed Recall@K, TES is monotonically decreasing in $|\text{retrieved}|$, with marginal penalty $\partial \text{TES}/\partial n = -\text{Recall}@K / (n(1+n) \cdot \ln(1+n)^{-1}...)$ capturing sub-linear cost growth.
\item \textbf{Boundedness}: $0 \leq \text{TES} \leq 1/\ln 2 \approx 1.44$ when $|\text{retrieved}| = 1$ and Recall@K = 1.0.
\end{enumerate}
These properties ensure TES penalizes context bloat proportionally to the marginal benefit of additional files.

\textbf{Empirical Validation.} We compute NDCG@5 for all 7 strategies across all 4 datasets (28 strategy-dataset pairs) and measure the Pearson correlation with TES. The overall correlation is $r = 0.87$ ($t = 9.05$, $\text{df} = 26$, $p < 0.001$), confirming that TES is a cost-adjusted quality metric.

\begin{table}[t]
\centering
\small
\caption{Comparison of retrieval metrics.}
\label{tab:metrics}
\begin{tabular}{@{}lccc@{}}
\toprule
\textbf{Metric} & \textbf{Measures} & \textbf{Cost-Aware?} \\
\midrule
Precision@K & Relevant fraction & No \\
Recall@K & Coverage of relevant items & No \\
NDCG@K & Ranking quality & No \\
F1@K & P/R balance & No \\
\textbf{TES} & \textbf{Recall-efficiency trade-off} & \textbf{Yes} \\
\bottomrule
\end{tabular}
\end{table}

\section{Experiments}

\subsection{Datasets}

\textbf{Synthetic Benchmark.} We generated a controlled benchmark with 50 source files and 166 queries. File references follow a Zipf distribution ($s = 1.2$). Queries are distributed across four trigger types: \textsc{Explicit\_Symbol} (79), \textsc{Semantic\_Concept} (47), \textsc{Temporal\_History} (10), and \textsc{Implicit\_Context} (30).

\textbf{Real Codebases.} We evaluated on three real Python projects:

\begin{table}[t]
\centering
\small
\caption{Dataset summary.}
\label{tab:datasets}
\begin{tabular}{@{}lcrrc@{}}
\toprule
\textbf{Dataset} & \textbf{Type} & \textbf{Files} & \textbf{Queries} & \textbf{Domain} \\
\midrule
Synthetic & Generated & 50 & 166 & Controlled \\
GraphPrompt & Real & 73 & 80 & Graph prompts \\
OneViral & Real & 299 & 84 & Social media \\
AgentNode & Real & 596 & 85 & Multi-agent \\
\bottomrule
\end{tabular}
\end{table}

Total: 968 real files across 3 codebases, 415 total queries.

\subsection{Baselines}

We compare CTX against seven baselines spanning the spectrum from no retrieval to production-grade neural RAG, graph-based retrieval, and a hybrid approach:

\begin{table}[t]
\centering
\small
\caption{Baseline strategies.}
\label{tab:baselines}
\begin{tabular}{@{}lp{5.5cm}@{}}
\toprule
\textbf{Strategy} & \textbf{Description} \\
\midrule
Full Context & Loads all files into context \\
BM25 & Sparse keyword retrieval \\
Dense TF-IDF & TF-IDF + cosine similarity \\
GraphRAG-lite & Import graph BFS without triggers \\
LlamaIndex & AST-aware code chunking + TF-IDF \\
Chroma Dense & ChromaDB + MiniLM-L6-v2 embeddings \\
Hybrid Dense+CTX & Dense seed + import graph expansion \\
CTX (Ours) & Trigger classification + graph traversal \\
\bottomrule
\end{tabular}
\end{table}

\subsection{Main Results}

\subsubsection{Synthetic Dataset (50 files, 166 queries)}

\begin{table}[t]
\centering
\small
\caption{Synthetic benchmark results.}
\label{tab:synthetic}
\begin{tabular}{@{}lccccc@{}}
\toprule
\textbf{Strategy} & \textbf{R@1} & \textbf{R@5} & \textbf{R@10} & \textbf{Tok\%} & \textbf{TES} \\
\midrule
Full Context & 0.014 & 0.075 & 0.170 & 100.0 & 0.019 \\
BM25 & 0.745 & 0.982 & 0.985 & 18.7 & 0.410 \\
Dense TF-IDF & 0.510 & 0.973 & 0.985 & 21.0 & 0.406 \\
GraphRAG-lite & 0.318 & 0.523 & 0.633 & 24.0 & 0.218 \\
LlamaIndex & 0.502 & 0.972 & 0.985 & 20.1 & 0.405 \\
Chroma Dense & 0.392 & 0.829 & 0.898 & 19.3 & 0.346 \\
Hybrid Dense+CTX & 0.392 & 0.725 & 0.757 & 23.6 & 0.303 \\
\textbf{CTX (Ours)} & \textbf{0.511} & \textbf{0.874} & \textbf{0.874} & \textbf{5.2} & \textbf{0.776} \\
\bottomrule
\end{tabular}
\end{table}

CTX achieves a TES of 0.776, which is 1.9$\times$ higher than BM25 (0.410) and 2.6$\times$ higher than Hybrid Dense+CTX (0.303). While BM25 achieves higher absolute Recall@5 (0.982 vs.\ 0.874), it consumes 3.6$\times$ more tokens.

\subsubsection{Real Codebases (3 projects, 968 files, 249 queries)}

\begin{table}[t]
\centering
\small
\caption{Real codebase results (average across 3 projects).}
\label{tab:real}
\begin{tabular}{@{}lcccc@{}}
\toprule
\textbf{Strategy} & \textbf{GP R@5} & \textbf{OV R@5} & \textbf{AN R@5} & \textbf{Avg TES} \\
\midrule
Full Context & 0.105 & 0.002 & 0.012 & 0.009 \\
BM25 & 0.457 & 0.156 & 0.252 & 0.116 \\
Dense TF-IDF & 0.531 & 0.224 & 0.300 & 0.146 \\
GraphRAG-lite & 0.567 & 0.206 & 0.066 & 0.107 \\
LlamaIndex & 0.545 & 0.366 & 0.188 & 0.145 \\
Chroma Dense & 0.458 & 0.252 & 0.162 & 0.118 \\
Hybrid Dense+CTX & 0.427 & --- & --- & --- \\
\textbf{CTX (Ours)} & 0.158 & 0.183 & 0.171 & \textbf{0.195} \\
\bottomrule
\end{tabular}
\end{table}

On real data, CTX achieves the highest average TES (0.195) using only 1.1\% of tokens.

\subsection{Trigger-Type Analysis}

\begin{table}[t]
\centering
\small
\caption{Synthetic Recall@5 by trigger type.}
\label{tab:trigger}
\begin{tabular}{@{}lcccc@{}}
\toprule
\textbf{Strategy} & \textbf{EXPL} & \textbf{SEMA} & \textbf{TEMP} & \textbf{IMPL} \\
\midrule
BM25 & 1.00 & 1.00 & 1.00 & 0.40 \\
Dense TF-IDF & 1.00 & 0.98 & 1.00 & 0.40 \\
Chroma Dense & 0.86 & 0.81 & 0.90 & 0.38 \\
LlamaIndex & 1.00 & 0.98 & 1.00 & 0.40 \\
Hybrid Dense+CTX & 0.81 & 0.66 & 0.60 & 0.53 \\
\textbf{CTX (Ours)} & \textbf{0.99} & 0.72 & \textbf{1.00} & \textbf{1.00} \\
\bottomrule
\end{tabular}
\end{table}

The critical distinction emerges on \textsc{Implicit\_Context} queries: CTX achieves perfect recall (1.0), while all baselines plateau at 0.33--0.43. The Hybrid variant achieves 0.53, improving over pure dense methods but below CTX due to imperfect seed selection.

\subsection{COIR-Style External Benchmark}

To position CTX against the broader code retrieval literature, we evaluate on a COIR-style benchmark constructed from the CodeSearchNet Python test set \citep{husain2019codesearchnet}.

\begin{table}[t]
\centering
\small
\caption{COIR-style evaluation (CodeSearchNet Python, 100 queries, 1000 corpus).}
\label{tab:coir}
\begin{tabular}{@{}lccc@{}}
\toprule
\textbf{Strategy} & \textbf{R@1} & \textbf{R@5} & \textbf{MRR} \\
\midrule
BM25 & 0.920 & 0.980 & 0.946 \\
Dense TF-IDF & 0.890 & 0.970 & 0.924 \\
Dense Embedding (MiniLM) & \textbf{0.960} & \textbf{1.000} & \textbf{0.978} \\
CTX Adaptive Trigger & 0.210 & 0.380 & 0.293 \\
Hybrid Dense+CTX & 0.930 & 0.950 & 0.940 \\
\bottomrule
\end{tabular}
\end{table}

On text-to-code retrieval, the Hybrid Dense+CTX variant achieves R@5 of 0.950, a 150\% improvement over CTX alone (0.380), nearly matching Dense Embedding (1.000). This confirms that the dense seed selection stage successfully bridges CTX's semantic gap, while the graph expansion provides additional structural context.

\subsection{LLM Downstream Quality}

We validate retrieval quality with an end-to-end code generation experiment using MiniMax M2.5 on 49 functions from the GraphPrompt codebase.

\begin{table}[t]
\centering
\small
\caption{LLM downstream quality (pass@1) on GraphPrompt ($n=49$).}
\label{tab:pass1}
\begin{tabular}{@{}lcccc@{}}
\toprule
\textbf{Strategy} & \textbf{pass@1} & \textbf{95\% CI} & \textbf{Tokens} & \textbf{Reduction} \\
\midrule
Full Context & 0.102 & [0.044, 0.218] & 11,952 & --- \\
CTX Adaptive & \textbf{0.265} & [0.162, 0.403] & 1,406 & 88.2\% \\
\bottomrule
\end{tabular}
\end{table}

CTX achieves 160\% higher pass@1 than Full Context (0.265 vs.\ 0.102) while using only 11.8\% of context tokens.

\section{Analysis}

\subsection{Why Import Graph Traversal Matters}

The \textsc{Implicit\_Context} results provide the strongest evidence for CTX's core thesis:

\begin{table}[t]
\centering
\small
\caption{\textsc{Implicit\_Context} Recall@5 comparison.}
\label{tab:implicit}
\begin{tabular}{@{}lcc@{}}
\toprule
\textbf{Approach} & \textbf{Best Strategy} & \textbf{R@5} \\
\midrule
Neural embedding & Chroma Dense & 0.33 \\
AST-aware chunking & LlamaIndex & 0.40 \\
Sparse keyword & BM25 / TF-IDF & 0.40 \\
Graph (no triggers) & GraphRAG-lite & 0.43 \\
Dense + graph & Hybrid Dense+CTX & 0.53 \\
Graph + triggers & \textbf{CTX} & \textbf{1.00} \\
\bottomrule
\end{tabular}
\end{table}

Only trigger-classified import graph traversal achieves perfect dependency resolution, with a margin of +133\% over the best baseline.

\subsection{Ablation Study}
\label{sec:ablation}

\begin{table}[t]
\centering
\small
\caption{Ablation study results.}
\label{tab:ablation}
\begin{tabular}{@{}llcc@{}}
\toprule
\textbf{Variant} & \textbf{Description} & \textbf{Synth R@5/TES} & \textbf{Real R@5/TES} \\
\midrule
Full CTX (A) & All components & 0.880/0.780 & 0.174/0.201 \\
No Graph (B) & Remove graph & 0.861/0.748 & 0.271/0.262 \\
No Classifier (C) & Remove triggers & 0.953/0.406 & 0.343/0.192 \\
Fixed-$k$=5 (D) & Remove adaptive-$k$ & 0.880/0.783 & 0.174/0.201 \\
\bottomrule
\end{tabular}
\end{table}

\textbf{Import graph (A vs B):} Removing graph traversal drops \textsc{Implicit\_Context} recall from 1.0 to 0.4 on synthetic data ($-60\%$).

\textbf{Trigger classifier (A vs C):} Without classification, the system achieves higher raw Recall@5 (0.953) but nearly doubles token usage, halving TES from 0.780 to 0.406.

\textbf{Adaptive-$k$ (A vs D):} Fixed $k=5$ performs nearly identically to full CTX, indicating the main efficiency gains come from trigger classification routing, not $k$-selection.

\subsection{Hybrid Dense+CTX Analysis}

The hybrid approach validates the complementary hypothesis:

\begin{itemize}
\item \textbf{COIR (text-to-code):} Hybrid R@5=0.950 vs CTX R@5=0.380 ($+150\%$). Dense seed selection bridges the semantic gap.
\item \textbf{IMPLICIT\_CONTEXT (synthetic):} Hybrid R@5=0.533 vs Chroma Dense R@5=0.383 ($+39\%$). Graph expansion adds structural context beyond what embeddings capture.
\item \textbf{Trade-off:} Hybrid uses 23.6\% tokens vs CTX's 5.2\%, sacrificing token efficiency for semantic coverage.
\item \textbf{Real codebases:} Hybrid achieves IMPLICIT\_CONTEXT R@5=0.229, outperforming CTX (0.076), because dense seed selection provides better initial file identification than rule-based matching on natural-language queries.
\end{itemize}

The optimal strategy depends on the deployment scenario: CTX for token-constrained dependency-heavy workloads; Hybrid for balanced workloads requiring both semantic and structural retrieval.

\subsection{Token Efficiency Frontier}

On synthetic data, sacrificing 10\% relative recall (0.982 to 0.874) yields a 72\% reduction in token usage. For cost-sensitive deployments processing 1,000 queries daily at 5.2\% token usage versus 18.7\%, CTX achieves a 3.6$\times$ cost reduction with minimal accuracy loss.

\section{Conclusion}

We presented CTX, a trigger-driven dynamic context loading system for code-aware LLM agents. CTX classifies developer queries into four trigger types and routes each to a specialized retrieval pipeline, with dependency-sensitive queries resolved through import graph traversal.

On a synthetic benchmark (50 files, 166 queries), CTX achieves a Trade-off Efficiency Score 1.9$\times$ higher than the best baseline while consuming only 5.2\% of total tokens. On three real Python codebases (968 files, 249 queries), CTX achieves the highest average TES (0.195) using only 1.1\% of tokens. On \textsc{Implicit\_Context} queries, CTX attains perfect Recall@5 of 1.0, compared to 0.4 for text-based baselines.

The Hybrid Dense+CTX variant demonstrates that dense retrieval and graph-based expansion are complementary: the hybrid achieves Recall@5 of 0.950 on the COIR benchmark (150\% improvement over CTX alone) while retaining structural dependency awareness (IMPLICIT\_CONTEXT R@5=0.533 vs 0.383 for dense-only methods). This validates the two-stage pipeline as a practical approach for balanced workloads.

\textbf{Limitations.} The pass@1 evaluation uses LLM self-judgment rather than execution-based verification, and the sample size ($n=49$) remains limited. The trigger classifier is rule-based; a learned classifier may improve robustness. On real codebases, CTX's absolute recall is lower than text-based baselines due to indexing pipeline limitations.

\textbf{Future Work.} We plan to: (1) expand evaluation with execution-based verification and larger sample sizes, (2) test on additional repositories in multiple languages, (3) replace the rule-based trigger classifier with a learned model, and (4) integrate \texttt{ast}-based parsing into the full CTX pipeline.

\bibliography{references}
\bibliographystyle{acl_natbib}

\end{document}
